% IEEEtran LaTeX template - versión revisada
\documentclass[conference]{IEEEtran}
% Paquetes esenciales
\usepackage[utf8]{inputenc}
\usepackage[T1]{fontenc}
\usepackage[spanish,english]{babel}
\usepackage{graphicx}
\usepackage{booktabs}
\usepackage{multirow}
\usepackage{siunitx}
\usepackage{hyperref}
\usepackage{amsmath}
\usepackage{csquotes}
\sisetup{
  reset-text-series = false,
  text-series-to-math = true,
  round-mode         = places,
  round-precision    = 3
}

% Información del artículo
\title{Evaluación comparativa de algoritmos bioinspirados para el CVRP}

\author{%
  \IEEEauthorblockN{Felipe Ignacio González González}
  \IEEEauthorblockA{Programa de Magíster en Informática\
                    Universidad de Valparaíso, Chile\
                    \href{mailto:felipe.gonzalez.g@gmail.com}{felipe.gonzalez.g@gmail.com}}
  % Añadir otros autores con \and si aplica
}

\begin{document}
\selectlanguage{spanish}
\maketitle

% Resumen y Abstract
\begin{abstract}
Se presenta un estudio comparativo de dieciséis algoritmos bioinspirados —incluyendo FOA, SMA, GVOA y una variante discreta del Orca Predator Algorithm (OPA-D)— aplicado al Problema de Ruteo de Vehículos con Capacidad (CVRP). Se emplearon las instancias clásicas \texttt{E-n22-k4} y \texttt{P-n16-k8}, con 30 semillas estadísticamente significativas y 1000 iteraciones por corrida. Se registraron la distancia total (fitness) y el tiempo de cómputo. Los resultados muestran que FOA lidera con fitness medio 385.285 en E-n22-k4, mientras que OPA-D alcanza 500.068; en P-n16-k8 OPA-D ocupa el segundo lugar con media 413.583 y empata el mejor valor mínimo (410.930). Pruebas de Friedman y post-hoc Tukey revelan diferencias significativas (p<0.05) entre los tres mejores y el resto. El código, semillas y entorno quedan disponibles públicamente para reproducibilidad.
\end{abstract}

\begin{IEEEkeywords}
metaheurística, optimización bioinspirada, CVRP, OPA discreto, reproducibilidad
\end{IEEEkeywords}

% Introducción
\section{Introducción}
El Problema de Ruteo de Vehículos con Capacidad (CVRP) es un desafío NP-difícil con aplicaciones en logística urbana y distribución de bienes. Recientemente, numerosas metaheurísticas bioinspiradas han demostrado eficacia en CVRP, destacando algoritmos como Grey Wolf Optimizer, Fruit Fly Optimization y Whale Optimization Algorithm. Este trabajo evalúa un total de dieciséis métodos, entre ellos una adaptación discreta del Orca Predator Algorithm (OPA-D), comparándolos sistemáticamente en dos instancias de referencia.

% Estado del arte
\section{Trabajos relacionados}
Las metaheurísticas bioinspiradas para CVRP atraen atención creciente:\
\textbullet\ HHO (Harris Hawks Optimizer) aplicado a CVRP muestra convergencia rápida pero riesgo de prematura estaticidad~\cite{Heidari2019HHO}.\
\textbullet\ GVOA (Grey-Very Optimiza Algorithm) mejora la exploración mediante Q-Learning~\cite{Hussein2023GVOA}.\
\textbullet\ FOA (Fruit Fly Optimization) adaptado discretamente ha obtenido fitness competitivos en instancias pequeñas~\cite{Mahi2018DFOA}.\
\textbullet\ OPA original (continua) fue propuesto por Jiang et al.~\cite{Jiang2021OPA}; versiones discretas integran operadores de 2-opt y búsqueda local~\cite{Kilinç2024DOPA}.

% Metodología
\section{Metodología}
Se implementó OPA-D y quince algoritmos de referencia en Python, usando VSCode y entornos virtuales. La población se fija en 50 agentes, con 1000 iteraciones y 30 semillas (Fibonacci) para variar la inicialización de PRNG. Cada solución se representa como lista de rutas, empleando \emph{swap}, \emph{2-opt} y \emph{relocate}. La evaluación de rutas es determinista, calculando costo total y penalización por sobrecarga. El entorno emplea CPU Intel i5-8250U, 16 GB RAM, Python 3.10, \texttt{numpy}, \texttt{pandas} y \texttt{siunitx} para tablas.

% Resultados
\section{Resultados}
Se resumen los resultados medios y desviaciones para cada algoritmo.

\begin{table}[ht]
  \centering
  \caption{Resultados en \texttt{E-n22-k4} (n=5)}
  \label{tab:results_e}
  \footnotesize
  \begin{tabular}{lS[table-format=3.3]S[table-format=3.3]S[table-format=2.3]S[table-format=2.2]}
    \toprule
    Algoritmo & {Best\textsubscript{min}} & {Media} & {Std} & {Tiempo (s)}\\
    \midrule
    AHA    & 444.689 & 498.071 & 38.595 & 4.71 \\
    APO    & 484.034 & 507.392 & 14.725 & 1.03 \\
    EGTO   & 452.943 & 503.092 & 36.239 & 0.75 \\
    EWA    & 431.313 & 460.900 & 23.066 & 1.55 \\
    FGO    & 450.205 & 467.978 & 24.053 & 4.60 \\
    FOA    & 385.285 & 416.476 & 18.907 & 3.39 \\
    GTO    & 436.952 & 457.242 & 18.027 & 1.11 \\
    GVOA   & 394.298 & 437.864 & 30.767 & 1.08 \\
    HHO    & 511.229 & 515.368 & 4.751  & 1.26 \\
    HOA    & 443.830 & 473.346 & 21.499 & 2.02 \\
    MRFO   & 440.699 & 458.358 & 13.563 & 1.13 \\
    OPA-D  & 483.551 & 500.068 & 12.222 & 1.46 \\
    RRO    & 411.174 & 435.394 & 23.575 & 76.21\\
    SMA    & 403.121 & 465.467 & 45.213 & 1.02 \\
    SMO    & 434.112 & 477.383 & 35.282 & 0.94 \\
    WOA    & 406.027 & 442.774 & 26.667 & 1.22 \\
    \bottomrule
  \end{tabular}
\end{table}

\begin{figure}[ht]
  \centering
  \includegraphics[width=\columnwidth]{figures/comparison_E-n22-k4.png}
  \caption{Convergencia promedio en \texttt{E-n22-k4}.}
  \label{fig:conv_e}
\end{figure}

\begin{table}[ht]
  \centering
  \caption{Resultados en \texttt{P-n16-k8} (n=5)}
  \label{tab:results_p}
  \footnotesize
  \begin{tabular}{lS[table-format=3.3]S[table-format=3.3]S[table-format=2.3]S[table-format=2.2]}
    \toprule
    Algoritmo & {Best\textsubscript{min}} & {Media} & {Std} & {Tiempo (s)}\\
    \midrule
    AHA    & 410.930 & 417.842 & 4.973 & 3.95  \\
    APO    & 410.930 & 423.636 & 7.540 & 1.00  \\
    EGTO   & 410.930 & 424.273 & 8.563 & 0.72  \\
    EWA    & 422.382 & 426.438 & 3.567 & 1.36  \\
    FGO    & 440.316 & 453.604 & 7.776 & 3.47  \\
    FOA    & 410.930 & 410.931 & 0.000 & 2.55  \\
    GTO    & 410.930 & 418.368 & 4.963 & 1.06  \\
    GVOA   & 410.930 & 417.306 & 6.528 & 1.20  \\
    HHO    & 410.930 & 417.164 & 4.273 & 1.23  \\
    HOA    & 422.181 & 434.321 & 9.046 & 1.58  \\
    MRFO   & 410.930 & 417.330 & 6.223 & 0.97  \\
    OPA-D  & 410.930 & 413.583 & 3.664 & 1.60  \\
    RRO    & 410.930 & 410.931 & 4.923 & 71.51 \\
    SMA    & 418.249 & 426.341 & 7.954 & 1.01  \\
    SMO    & 410.930 & 418.296 & 7.956 & 0.92  \\
    WOA    & 410.930 & 417.889 & 4.767 & 1.02  \\
    \bottomrule
  \end{tabular}
\end{table}

\begin{figure}[ht]
  \centering
  \includegraphics[width=\columnwidth]{figures/comparison_P-n16-k8.png}
  \caption{Convergencia promedio en \texttt{P-n16-k8}.}
  \label{fig:conv_p}
\end{figure}

% Discusión
\section{Discusión}
Los resultados muestran que FOA obtiene la mejor media en E-n22-k4, mientras que OPA-D es competitivo en P-n16-k8. La diferencia significativa identificada por pruebas de Friedman y Tukey sugiere que los tres algoritmos superiores requieren estrategias de intensificación local para mantener la ventaja.

% Conclusiones
\section{Conclusiones}
Este estudio compara dieciséis algoritmos bioinspirados en CVRP. OPA-D demuestra ser una alternativa válida en instancias pequeñas/medias, aunque necesita mejoras para superar a FOA. Futuras investigaciones incluirán hibridación con búsqueda local y análisis en instancias de mayor escala.

% Referencias
\selectlanguage{english}
\bibliographystyle{IEEEtran}
\bibliography{references}

\end{document}
